\documentclass
% [aspectratio=169]
{beamer}

\usepackage{CJKutf8}
\usepackage{listings}
\usepackage{tikz}
\usepackage{hyperref}
\usepackage{xcolor}
\usepackage{verbatim}
\usepackage{eso-pic}
\usepackage{courier}
\usepackage{textcomp}

\newcommand{\link}[1]{\href{#1}{#1}}

\newcommand{\insertGraph}[3]{

	\centerline{\includegraphics[scale=#1]{#2}} % 0/0

	\centerline{#3}

}

\newcommand{\HCCLogoSimp}{
	\begin{tikzpicture}[scale=0.1]
		\definecolor{_00ccff}{HTML}{00ccff}
\def\center{(54.9655,-19.629)}
\def\radius{7.924}
\fill [color=_00ccff]
	(55.371,-11.075) -- (57.286,-11.417) -- (55.693,-20.323) -- 
	(53.803,-19.976);
\begin{scope}
	% (47.042,-11.705) rectangle (62.889,-27.553)
	% (49.018,-13.682) rectangle (60.913,-25.576)
%	\clip
%		(59.4605,-11.563) -- (54.9665,-19.629) -- (53.3695,-11.511) -- 
%		(47.041,-11.511) -- (47.041,-27.553) -- (62.890,-27.553) -- 
%		(62.890,-11.511);
%	\fill [color=_00ccff, even odd rule]
%		\center circle [radius=5.9475]
%		(59.4605,-11.563) -- (53.3695,-11.511) -- (54.9665,-19.629)
%		\center circle [radius=\radius]; 
	% https://tex.stackexchange.com/questions/510281/tikz-fill-only-the-a-b-c
	\path \center circle [radius=\radius];
	\clip[overlay]
		(53.3695,-11.511) -- (59.4605,-11.563) -- (54.9665,-19.629) [rev];
	\clip[overlay, eo] \center circle [radius=5.9475, rev];
	\fill[color=_00ccff] \center circle [radius=\radius];
\end{scope}


	\end{tikzpicture}
}

\newcommand{\HCCLogoFull}{
	\begin{tikzpicture}[scale=0.05]
		\definecolor{_00ccff}{HTML}{00ccff}
\def\center{(54.9655,-19.629)}
\def\radius{7.924}
\fill [color=_00ccff]
	(55.371,-11.075) -- (57.286,-11.417) -- (55.693,-20.323) -- 
	(53.803,-19.976);
\begin{scope}
	% (47.042,-11.705) rectangle (62.889,-27.553)
	% (49.018,-13.682) rectangle (60.913,-25.576)
%	\clip
%		(59.4605,-11.563) -- (54.9665,-19.629) -- (53.3695,-11.511) -- 
%		(47.041,-11.511) -- (47.041,-27.553) -- (62.890,-27.553) -- 
%		(62.890,-11.511);
%	\fill [color=_00ccff, even odd rule]
%		\center circle [radius=5.9475]
%		(59.4605,-11.563) -- (53.3695,-11.511) -- (54.9665,-19.629)
%		\center circle [radius=\radius]; 
	% https://tex.stackexchange.com/questions/510281/tikz-fill-only-the-a-b-c
	\path \center circle [radius=\radius];
	\clip[overlay]
		(53.3695,-11.511) -- (59.4605,-11.563) -- (54.9665,-19.629) [rev];
	\clip[overlay, eo] \center circle [radius=5.9475, rev];
	\fill[color=_00ccff] \center circle [radius=\radius];
\end{scope}


		\definecolor{_004455}{HTML}{004455}
\definecolor{_006680}{HTML}{006680}
\definecolor{_0088aa}{HTML}{0088aa}
\definecolor{_00aad4}{HTML}{00aad4}
\fill [color=_00aad4] (49.0,-46.0) rectangle (64.0,-50.5);
\fill [color=_0088aa] (44.5,-35.5) rectangle (49.0,-50.5);
\fill [color=_006680] (44.5,-31.0) rectangle (64.0,-35.5);
\fill [color=_00aad4] (27.5,-46.0) rectangle (42.5,-50.5);
\fill [color=_0088aa] (22.5,-35.5) rectangle (27.5,-50.5);
\fill [color=_006680] (22.5,-31.0) rectangle (42.0,-35.5);
\fill [color=_0088aa] (16.0,-35.5) rectangle (20.5,-51.0);
\fill [color=_006680] (4.5,-31.0) rectangle (20.5,-35.5);
\fill [color=_004455] (0.0,-24.0) rectangle (4.5,-53.5);
% \draw [color=red] (0.0,0.0) rectangle (64,-64);

	\end{tikzpicture}
}

\setbeamercolor{background canvas}{bg=}

\newcommand{\PreFirstFrame}{
	\AddToShipoutPictureFG*{
		\AtPageLowerLeft{
			\put(\LenToUnit{0.05\paperwidth},\LenToUnit{0.1\paperheight}){
				\footnotesize
				这个指引文档在
				\href{https://creativecommons.org/licenses/by-sa/3.0/deed.zh}
				{知识共享 署名-相同方式共享 3.0协议}之条款下提供
			}
			\put(\LenToUnit{0.05\paperwidth},\LenToUnit{0.05\paperheight}){
				\footnotesize
				This guidance is available under the 
				\href{https://creativecommons.org/licenses/by-sa/3.0/}
				{Creative Commons Attribution-ShareAlike License}
			}
			\put(\LenToUnit{0.6\paperwidth},\LenToUnit{0.15\paperheight}){
				\HCCLogoFull
			}
		}
	}
}

\newcommand{\PostFirstFrame}{
	\AddToShipoutPictureBG{
		\AtPageLowerLeft{
			\put(\LenToUnit{0.8\paperwidth},\LenToUnit{0.15\paperheight}){
				\HCCLogoSimp
			}
		}
	}
}

\newcommand{\PreLastFrame}{
	\ClearShipoutPictureBG

	\AddToShipoutPictureFG*{
		\AtPageLowerLeft{
			\put(\LenToUnit{0.6\paperwidth},\LenToUnit{0.15\paperheight}){
				\HCCLogoFull
			}
		}
	}
}

% Note: this C style differs a lot from gedit's
\lstdefinestyle{cstyle}{
	language=c,
	basicstyle=\ttfamily,
	morekeywords={with},
	keywordstyle=\bfseries\color[HTML]{a52a2a},	
	commentstyle=\color[HTML]{0000ff},
	stringstyle=\color[HTML]{ff0bff},
	keywordstyle=[3]\color[HTML]{008a8c},
	alsoletter={0,1,2,3,4,5,6,7,8,9,.},
	morekeywords=[4]{0,1,2,100,999},
	keywordstyle=[4]\color[HTML]{ff0bff},
	upquote=true,
	breaklines=true,
}

\lstdefinestyle{pythonstyle}{
	language=python,
	basicstyle=\ttfamily,
	% frame=single,
	morekeywords={with,yield},
	keywordstyle=\bfseries\color[HTML]{a52a2a},	
	keywordstyle=[2]\color[HTML]{008a8c},
	commentstyle=\color[HTML]{0000ff},
	stringstyle=\color[HTML]{ff0bff},
	keywordstyle=[3]\color[HTML]{008a8c},
	alsoletter={0123456789.},
	morekeywords=[4]{False,True,
		0,1,2,3,4,5,6,7,8,9,10,11,12,13,15,17,19,16,20,24,27,31,32,33,34,35,38,
		45,56,60,64,81,95,97,99,100,123,243,256,400,512,576,729,999,1024,1234,
		1365,1366,2000,2187,2836,2957,3856,3857,5678,6561,9274,100000,1000000,
		0.5,3.14,3.4,
		0x1234,},
	keywordstyle=[4]\color[HTML]{ff0bff},
	upquote=true,
	breaklines=true,
	showstringspaces=false,
}

\lstset{
	tabsize=4,
	columns=fixed,
	extendedchars=false,
}

\newcommand{\inlinePython}{\lstinline[style=pythonstyle]}



\begin{document}
\begin{CJK}{UTF8}{gbsn}

\PreFirstFrame
\begin{frame} [fragile]
	\centerline{\fontsize{42}{42}\selectfont Python Talk 8}
\end{frame}
\PostFirstFrame

\begin{frame} [fragile]
	\frametitle{连线：这些功能可以做什么}
	\linespread{1.5}
	\begin{columns}[T]
		\begin{column}[T]{.5\textwidth}
			\begin{enumerate}
			\item 产生随机数
			\item 制造唯一链接
			\item 安全输入密码
			\item 生成二维码
			\item 架设一个网站
			\end{enumerate}
		\end{column}
		\begin{column}[T]{.5\textwidth}
			\begin{itemize}
			\item \inlinePython{uuid}
			\item \inlinePython{qrcode}
			\item \inlinePython{python3}
			\item \inlinePython{getch}
			\item \inlinePython{random}
			\end{itemize}
		\end{column}
	\end{columns}
\end{frame}

\begin{frame} [fragile]
	\frametitle{实际应用}
	\linespread{1.5}
	\begin{itemize}
	\item 在这次讲座中我们将会用Python完成一些实际的工作
		\begin{itemize}
		\item 生成二维码
		\item 批量处理文本
		\item 裁剪图片
		\item 批量下载附件
		\end{itemize}
	\end{itemize}
\end{frame}

\begin{frame} [fragile]
	\frametitle{\inlinePython{qrcode}}
	\linespread{1.5}
	\begin{itemize}
	\item 回忆：\inlinePython{qrcode}用来生成二维码
		\begin{lstlisting}[style=pythonstyle, gobble=8, texcl, escapechar=@]
		import qrcode
		qrcode.run_example('@\color[HTML]{ff0bff}内容字符串@')
		\end{lstlisting}
	\item 需要额外安装软件包：\texttt{pip3 install qrcode}
	\item 任务：生成十一圈活动二维码
		\begin{itemize}
		\item 例：\texttt{http://shiyiquan.net/event/123/}
		\item 将123替换为活动id（通过input得到）
		\end{itemize}
	\end{itemize}
\end{frame}

\begin{frame} [fragile]
	\frametitle{批量处理文本}
	\linespread{1.25}
	\begin{itemize}
	\item 按行读取文件
		\begin{lstlisting}[style=pythonstyle, gobble=8, texcl]
		for i in 文件.readlines() :
			...			 # i 是文件的其中一行
		\end{lstlisting}
	\item 取出数据
		\begin{lstlisting}[style=pythonstyle, gobble=8, texcl, escapechar=@]
		'@\color[HTML]{ff0bff}字符串@'.split()	# 用 [] 选择需要的部分
		\end{lstlisting}
	\item 任务
		\begin{enumerate}
		\item 在命令行执行：\texttt{ls -l /etc > /tmp/etc}
		\item 计算 \texttt{/tmp/etc} 文件第一行以外每一行的第五列的和
		\end{enumerate}
	\item 答案
		\fontsize{2}{2}\selectfont
		\begin{lstlisting}[style=pythonstyle, gobble=8, texcl]
		s = 0
		for i in open('/tmp/etc').readlines()[1:] :
			s += int(i.split()[4])

		print(s)
		\end{lstlisting}
	\end{itemize}
\end{frame}

\begin{frame} [fragile]
	\frametitle{图片处理}
	\begin{itemize}
	\item Pillow库
		\begin{itemize}
		\item \inlinePython{from PIL import Image}
		\item 需要额外安装软件包：\texttt{pip3 install pillow}
		\end{itemize}
	\begin{lstlisting}[style=pythonstyle, gobble=4, texcl, escapechar=@]
	a = Image.open('@\color[HTML]{ff0bff}文件名@')		# 打开图像
	a.show()					 # 预览
	b = a.crop((@左@, @上@, @右@, @下@))  # 裁剪
	b.save('@\color[HTML]{ff0bff}文件名@')				# 保存
	\end{lstlisting}
	\end{itemize}
	\begin{columns}[T]
		\begin{column}[T]{.5\textwidth}
			\begin{itemize}
			\item 任务
				\begin{itemize}
				\item 编写一个Python程序，将输入的任何图片从中间裁减为正方形
				\item 参考右侧示例，将深色部分从浅色的输入图片中裁减出来
				\end{itemize}
			\end{itemize}
		\end{column}
		\begin{column}[T]{.5\textwidth}
			\begin{tikzpicture}[scale=0.5]
			\def\x{7}
			\filldraw [color={rgb,255:red,255; green,255; blue,255}]
				(-0.5,-0.5) rectangle (\x+4.5,6.5);
			\filldraw [color={rgb,255:red,0; green,170; blue,212}]
				(0,1) rectangle (6,5);
			\filldraw [color={rgb,255:red,0; green,170; blue,212}]
				(\x,0) rectangle (\x+4,6);
			\filldraw [color={rgb,255:red,0; green,102; blue,128}]
				(1,1) rectangle (5,5);
			\filldraw [color={rgb,255:red,0; green,102; blue,128}]
				(\x,1) rectangle (\x+4,5);
			\end{tikzpicture}
		\end{column}
	\end{columns}
\end{frame}

\begin{frame} [fragile]
	\frametitle{批量下载附件}
	\begin{itemize}
	\item 生成命令来批量下载这个网页上的附件
		\begin{itemize}
		\item \link{http://c.hcc.io/f/20151002moon/}
		\item 提示：此网站有时可能无法访问，此时可随意找互联网上的网页尝试
		\end{itemize}
	\item 下载网页 / 附件的命令
		\begin{itemize}
		\item \texttt{curl 网址 > 文件}
		\item 用\inlinePython{os.system}运行
		\end{itemize}
	\item 以{\color{blue}蓝色}字符为参照，截取网页每行的信息
	\begin{lstlisting}[style=pythonstyle, gobble=4, texcl, escapechar=@]
	@\color{blue}<tbody>@
		@\color{blue}<tr>@...</tr>
		@\color{blue}<tr>@...</tr>
		...
	@\color{blue}</tbody>@
	\end{lstlisting}
	\end{itemize}
\end{frame}

\begin{frame} [fragile]
	\frametitle{从每一行中取得信息}
	\begin{itemize}
	\item 以{\color{blue}蓝色}字符为参照，截取{\color[HTML]{ff0bff}粉色}的内容信息
	\item 善用 \inlinePython{split}
	\end{itemize}
	\begin{lstlisting}[basicstyle=\ttfamily, gobble=4, texcl, escapechar=@]
	<tr>
	<th scope="row" class="hidden-xs"><h5>4</h5></th>
	<td>@\color{blue}<h5>\color[HTML]{ff0bff}月亮.MOV\color{blue}</h5>@</td>
	<td class="xs"><h5>459.6 MiB</h5></td>
	<td class="xs"><h5>2</h5></td>
	<td>@\color{blue}<a href="\color[HTML]{ff0bff}/f/20151002moon/\color{blue}"@ class="b">
	<span class="glyphicon"></span> Save</a>
	</td>
	</tr>
	\end{lstlisting}
\end{frame}

\begin{frame} [fragile]
	\frametitle{解法}
	\begin{lstlisting}[style=pythonstyle, gobble=4, texcl, escapechar=@]
	for i in a.split('<tbody>')[2] \
		.split('</tbody>')[0].split('<tr>')[1:] : 
		print('curl %s > %s' % 
			('http://c.hcc.io' + \
			i.split('href="')[1].split('"')[0], 
			i.split('<h5>')[2].split('</h5>')[0]))
	\end{lstlisting}
	\begin{itemize}
	\item 使用正则表达式可以让程序更加简洁。这将在以后介绍。
	\end{itemize}
\end{frame}

\PreLastFrame
\begin{frame}
	\centerline{\fontsize{32}{32}\selectfont 感谢参加此次活动}
\end{frame}

\newpage
\end{CJK}
\end{document}

